\documentclass[10pt,a4paper]{amsart}
\usepackage[margin=19mm]{geometry}
\usepackage{amsmath,amssymb,mathtools}
\usepackage[scr=rsfs,cal=euler]{mathalfa}
\usepackage{xfrac}
\usepackage[unicode,hidelinks,bookmarks=false]{hyperref}
\newcommand{\ie}{i.e.\ }
\newcommand{\cf}{cf.\ }
\newcommand{\resp}{respectively\ }
\newcommand{\Subsection}{\subsection}
\providecommand{\nobreakdash}{\nobreak\hbox{-}}
\setlength{\emergencystretch}{2em}
\multlinegap0pt
\usepackage{bm}
\usepackage{bbm}
\usepackage{dsfont}
\usepackage{xspace}
\usepackage{cancel}
\usepackage{mathtools}
\usepackage{amsthm}
\makeatletter
\@ifundefined{theorem}{\newtheorem{theorem}{Theorem}}{}
\makeatother
\DeclareMathOperator*{\argmax}{\text{arg\,max}}
\DeclareMathOperator*{\argmin}{\text{arg\,min}}
\newcommand{\lag}{\mathcal{L}}
\title[The relation of bias with risk in empirically constrained in]{The relation of bias with risk in empirically constrained inferences}
\author{Dalton A R Sakthivadivel}
\date{Source: arXiv:2507.03699v1 (CC BY 4.0)}
\begin{document}
\maketitle

\noindent\emph{Source and licence.} The relation of bias with risk in empirically constrained inferences. Version arXiv:2507.03699v1 (CC BY 4.0), distributed under the Creative Commons Attribution 4.0 International licence, as recorded in the arXiv licence metadata for this exact version and shown on the abstract page https://arxiv.org/abs/2507.03699v1. Full rights notice: licenses/arxiv-2507.03699-CC-BY-4.0.txt.\par\smallskip

\noindent\textit{Editorial scope.} Counted unit: the Theorem whose source label is \texttt{max-likelihood-thm} in the article (source0.tex lines 407--409, source label \texttt{max-likelihood-thm}), delivered together with the article's own complete original proof of it (source0.tex lines 410--436, one \texttt{proof} environment of 27 lines). No mathematics is written, repaired or re-derived: statement and proof are the article's own text line for line. Required context retained uncounted: level-three-thm (lines 383--392). Every definition, display and label of those lines is retained unchanged. Omitted from this package and not redistributed: 1--382, 393--406, 437--482. Nothing inside a retained excerpt is omitted or altered: every retained body is delivered exactly as the article's lines are written, so no omission receipt of any kind accompanies this package. Citations: the retained excerpts carry no citation command at all, so no bibliography entry of the article is transcribed and no reference is redistributed. Reference adaptation: none was needed and none was made; every reference inside the delivered unit resolves inside it, so no unresolved reference or citation is produced. Printed slips kept literal: no printed word, symbol, hypothesis or number of the counted statement or of its proof was altered. Layout and class: the article's own class is \texttt{amsart} with packages including inputenc, geometry, amssymb, amsmath, amsthm, amsfonts, graphicx, color, mathrsfs, physics, mathtools, tikz, faktor, tikz-cd; the wrapper is the lane's pinned \texttt{amsart} editorial wrapper at 10pt on A4, which restates the theorem-like environments the retained text needs and adds the article's own macro definitions that the retained text needs (\texttt{\textbackslash argmax}, \texttt{\textbackslash argmin}, \texttt{\textbackslash lag}), with extra wrapper packages (bm, bbm, dsfont, xspace, cancel, mathtools). The delivered numbering is the unit's own. Assets: no retained span contains a figure environment, a graphics-inclusion command, a PDF-page inclusion, a table or a TikZ picture; no image file is copied into this package, the package declares no asset dependency and no image file is redistributed with it. Licence: version arXiv:2507.03699v1 of this article is distributed under the Creative Commons Attribution 4.0 International licence, as recorded in the arXiv licence metadata for this exact version and shown on the abstract page https://arxiv.org/abs/2507.03699v1. A full rights notice is delivered at licenses/arxiv-2507.03699-CC-BY-4.0.txt. Characters: every retained excerpt is pure ASCII, so the delivered engine, XeLaTeX with its default Latin Modern text font, reports no missing character.
\begin{theorem}\label{level-three-thm}
    There exists a level three large deviations principle associated to the empirical conditional probabilities of empirical measures with prescribed means 
    \[
        M_k = P^n(\{\mu : \lag\cdot\mu \in \Xi\} \mid \mu \in K),
    \]
    with rate function 
    \[
        -\inf_{\nu \in H} D_{\mathrm{KL}}(\nu \| Q).
    \]
\end{theorem}

\begin{theorem}\label{max-likelihood-thm}
    Computing the most probable empirical measure given some range of constraint values is equivalent to a certain maximum likelihood problem.
\end{theorem}

\begin{proof}
    We will begin by recalling that
    \[
        P^n(\mu \in K \mid \lag\cdot\mu \in \Xi) \propto P^n(\{\mu : \lag\cdot\mu \in \Xi\} \mid \mu \in K) P^n(\mu \in K).
    \]
    Theorem \ref{level-three-thm} gives us an explicit expression for the optimal estimate of the first factor by putting
    \[
        \nu^* = \argmin_{\{\nu : U\nu = \eta\}} D_{\mathrm{KL}}\Big(\nu(\{\mu\}_\Xi \mid \mu \in K) \| Q(\{\mu\}_\Xi \mid \mu \in K)\Big).
    \]
    Solving the convex dual problem to obtain an expression for the minimiser yields (denoting the value of the satisfactory Lagrange multiplier as $\lambda_\eta$)
    \[
        \nu^* = e^{-\lambda_\eta U(\{\mu\}_\Xi)}Q(\{\mu\}_\Xi)
    \]
    from which we have 
    \begin{align*}
        \frac{1}{n}\log &P^n(\mu \in K \mid \lag \cdot \mu \in \Xi, U\cdot \nu = \eta) \\ &\asymp -\inf_{\mu \in K}\left(\inf_{\mu \in K \cap \{\mu\}_\Xi}D_{\mathrm{KL}}(\mu \| P) + %\inf_{\{\nu : U\nu = \eta\}} 
        \lambda_\eta U(\{\mu\}_\Xi) - \log Q(\{\mu\}_\Xi) \right).
    \end{align*}
    Using additivity and simplifying all three terms, we arrive at
    \[
        -\inf_{\mu \in K \cap \{\mu\}_\Xi} D_{\mathrm{KL}}(\mu \| P) -\inf_{\mu \in K \cap \{\mu\}_\Xi}\lambda_\eta U(\mu) + \sup_{\mu \in K \cap \{\mu\}_\Xi} Q(\mu).
    \]
    It follows that we have converted our study of the point of concentration to the easily soluble maximum likelihood estimate
    \begin{equation}\label{max-likelihood-eq}
        \argmax_{\mu \in K \cap \{\mu\}_\Xi} e^{- D_{\mathrm{KL}}(\mu \| P)} e^{-\lambda_\eta U(\mu)}Q(\mu).
    \end{equation}
\end{proof}

\end{document}
